\chapter{SpriteBatch and 2D Rendering}
\label{ch:spritebatch}

\cnaclass{SpriteBatch} is the API most CNA code touches most often, and --- because it is
also the one class all 46 public renderer identities in Part~\ref{part:renderers} must implement --- the class
whose bugs and fixes best illustrate what ``renderer parity'' actually costs to achieve. This
chapter covers its API surface, its supporting types (\cnaclass{SpriteSortMode},
\cnaclass{SpriteEffects}, \cnaclass{SpriteFont}), and the specific, real bugs found while
verifying it, because the bugs are as instructive as the API itself.

\section{Begin, Draw, End}

\cnaclass{SpriteBatch} inherits \cnaclass{GraphicsResource} and offers five \texttt{Begin()}
overloads of increasing specificity. The public contract defaults whatever the caller omits:
\texttt{Begin()} alone defaults to \texttt{SpriteSortMode::Deferred},
\cnaclass{BlendState::AlphaBlend}, a \texttt{LinearClamp} sampler, no depth-stencil state,
\texttt{RasterizerState::\allowbreak{}CullCounterClockwise}, and an identity transform; the fuller
overloads let a caller override the sort mode and blend state, the three optional state
objects (nullptr falls back to the same defaults), a custom \cnaclass{Effect} (nullptr falls
back to an internal default \cnaclass{SpriteEffect}), and finally an explicit transform
\cnaclass{Matrix} applied to every sprite in the batch before projection.

\begin{sourcenote}
\textbf{The live CNA implementation does not currently honor that whole state contract.}
\texttt{SpriteBatch::Begin()} applies the requested/default blend and depth-stencil values and
forwards the resolved sampler, but its \texttt{RasterizerState*} parameter is commented out and
never read. FNA stores the caller's value (or
\texttt{CullCounterClockwise}, the matching \cnaclass{RasterizerState} preset) and applies it in
\texttt{PrepRenderState()}. CNA instead leaves whatever rasterizer state was already on
\cnaclass{GraphicsDevice}. The practical scissor consequence is direct: passing a rasterizer
with \texttt{ScissorTestEnable=true} to the normal full \texttt{Begin()} overload cannot enable
clipping on any renderer; a previously-enabled state can also leak through a later default
\texttt{Begin()}. A game can work around the first half only by assigning
\texttt{GraphicsDevice.RasterizerState} separately before drawing. Section
\ref{sec:backend-viewport-scissor-matrix} then determines whether the selected renderer can
actually honor it.
\end{sourcenote}

\texttt{Draw} itself has nine overloads: six are the canonical XNA forms ---
position-or-destination-rectangle, an optional source rectangle (represented as
\texttt{std::optional<Rectangle>} rather than C\#'s nullable \texttt{Rectangle?}), color, and
progressively more control over rotation, origin, scale (as either a uniform \texttt{float} or
a non-uniform \cnaclass{Vector2}), \cnaclass{SpriteEffects}, and layer depth --- and three are
\texttt{CNAEXT} raw texture-and-rectangle convenience forms layered on top.
\texttt{DrawString} adds six more overloads (\texttt{string} or \texttt{StringBuilder}, each in
a simple 4-argument form, a uniform-scale 9-argument form, and a non-uniform-scale
(\cnaclass{Vector2} scale) 9-argument form).

\section{The Draw overload family, and a worked rotation example}

The nine \texttt{Draw} overloads split cleanly into two families by their second parameter ---
a \cnaclass{Vector2} \texttt{position} (four overloads, increasingly parameterized) or a
\cnaclass{Rectangle} \texttt{destinationRectangle} that scales the texture to fit (three
overloads) --- plus the two \texttt{CNAEXT} raw \texttt{float x, float y} convenience forms
already noted. Every real (non-\texttt{CNAEXT}) overload takes an explicit \texttt{Color} tint
--- there is no overload that omits it, unlike the \texttt{CNAEXT} convenience forms. The table
below abbreviates each overload's parameter names for width (\texttt{tex} for \texttt{texture},
\texttt{pos} for \texttt{position}, and so on) --- the full, exact parameter names for every
overload are given in the fully-parameterized worked example immediately below the table.

\begin{center}
\begin{longtable}{>{\raggedright\arraybackslash}p{0.30\textwidth}>{\raggedright\arraybackslash}p{0.58\textwidth}}
\toprule
\textbf{Family member} & \textbf{Adds, relative to the simplest form in its family} \\
\midrule
\endhead
\texttt{Draw(tex, pos, color)} & The simplest real-XNA form: whole texture, no rotation, unit
  scale. \\[4pt]
\texttt{Draw(tex, pos, srcRect, color)} & An optional source sub-rectangle
  (\texttt{std::optional<Rectangle>}, not C\#'s nullable \texttt{Rectangle?}) --- the
  sprite-sheet-cell pattern. \\[4pt]
\texttt{Draw(tex, pos, srcRect, color, rot, origin, scale, fx, depth)} & Uniform
  \texttt{float} scale. \\[4pt]
\texttt{Draw(tex, pos, srcRect, color, rot, origin, scale, fx, depth)} & Same, but
  \texttt{scale} is a non-uniform \cnaclass{Vector2} --- these two fully-parameterized
  overloads differ \emph{only} in that one field's type. \\[4pt]
\texttt{Draw(tex, destRect, color)} & Whole texture, scaled to fill \texttt{destRect}. \\[4pt]
\texttt{Draw(tex, destRect, srcRect, color)} & Adds the source sub-rectangle. \\[4pt]
\texttt{Draw(tex, destRect, srcRect, color, rot, origin, fx, depth)} & Full control,
  destination-rectangle family --- note there is no separate \texttt{scale} parameter here at
  all, since the destination rectangle already implies it. \\
\bottomrule
\end{longtable}
\end{center}

\texttt{origin} is the parameter every one of the fully-parameterized overloads shares, and it
is worth a worked example precisely because \S\ref{sec:sdl-renderer-two-bugs} below names a
real bug caused by getting a pivot convention wrong: \texttt{origin} is a point \emph{in
texture space} (not screen space) that both the rotation pivot and the drawn position are
measured relative to. Rotating a $64\times64$ sprite around its own center, rather than its
top-left corner, means passing that texture's own center as \texttt{origin}:

\begin{lstlisting}
Vector2 center(texture.getWidthProperty() / 2.0f, texture.getHeightProperty() / 2.0f);
float angle = totalSeconds; // radians; grows every frame for a continuous spin

spriteBatch_->Draw(
    texture, spritePosition, std::nullopt, Color::White,
    angle, center, 1.0f, SpriteEffects::None, 0.0f);
\end{lstlisting}

Passing \texttt{Vector2::Zero} for \texttt{origin} instead --- the easy mistake, since it is
also the type's default-constructed value --- rotates the sprite around its top-left corner,
which reads as the sprite ``orbiting'' \texttt{spritePosition} in a circle rather than spinning
in place; the visual symptom is easy to misdiagnose as a rotation-math bug when the actual
cause is simply the wrong \texttt{origin}.

\section{Vector positions are quantized before they reach a renderer}

The \cnaclass{Vector2}-position overloads do not preserve sub-pixel destinations. Their shared
queue record stores an integer \cnaclass{Rectangle}; CNA validates the position and scaled size,
then converts each floating component with truncation toward zero. Thus \texttt{(10.9f, 4.9f)}
and \texttt{(10.1f, 4.1f)} both enter every renderer as \texttt{(10,4)}. This is a shared
\cnaclass{SpriteBatch} fidelity difference from XNA, not a renderer-specific rasterization rule.

Text follows a subtly different policy. \texttt{DrawString()} computes each transformed glyph
destination in floating point and rounds it to the nearest integer before storing the same
rectangle shape. A sprite and a glyph submitted at the same fractional coordinate can therefore
land on adjacent pixels. Both conversion helpers reject non-finite values and results outside
the signed 32-bit range; the safety check prevents undefined conversion, but does not restore the
lost fractional position. Code that depends on smooth sub-pixel camera motion must account for
this quantization above the renderer layer.

\section{The D3D9 half-pixel offset: real, renderer-local, and a genuine mutation-testing trap}

Real XNA~4.0 ran exclusively on Direct3D~9, whose rasterizer places texel centers at integer
pixel coordinates rather than pixel centers --- the well-known ``D3D9 half-pixel offset''
problem, notorious enough that porting guides for other 2D APIs still warn about it decades
later. \cnaclass{SpriteBatch}'s own shared, renderer-agnostic \texttt{SpriteBatch.cpp} carries
no trace of this compensation at all, which is exactly correct behavior for every renderer built
on a modern API (EasyGL, Vulkan, BGFX, D3D11, D3D12, WebGPU): none of them share Direct3D~9's
own texel-center convention, so there is nothing to compensate for, and adding an offset
unconditionally would introduce a real bug on every renderer that does not need one. On this
one renderer where the convention genuinely applies, the compensation lives one level down,
inside \texttt{D3D9SpriteBatchRenderer::BuildMatrixTransformEXT} specifically --- baked directly
into the projection matrix \cnaclass{SpriteBatch} constructs, the same technique the classic
XNA~4.0 D3D9 fix itself used:

\begin{lstlisting}
// Worked example: the real D3D9-only half-pixel compensation, folded
// into the orthographic projection SpriteBatch's own MatrixTransform
// applies -- not a separate step, not visible anywhere in shared code.
Matrix projection = Matrix::CreateOrthographicOffCenter(0, width, height, 0, 0, 1);
projection.M41 += -0.5f * projection.M11;
projection.M42 += -0.5f * projection.M22;
// MatrixTransform = SpriteBatch's own transform, then this projection
// (row-vector convention: transform applied first).
\end{lstlisting}

Verifying this against the real XNA~4.0 oracle (Chapter~\ref{ch:direct3d-backends}) surfaced a
genuine mutation-testing trap worth knowing as a general lesson, not just a D3D9 specific:
\textbf{two entirely different test designs independently turned out to be structurally
incapable of detecting whether this offset was even present at all.} A $1\times1$ test texture
cannot detect it, because the half-pixel offset shifts \emph{which texel a given screen pixel
samples}, not where a sprite's geometric edges land on screen --- a texture with only one texel
has nothing else to shift onto, so removing the offset entirely changes nothing a
one-texel-texture test can observe. Separately, sampling only at a drawn rectangle's own
quadrant-center points misses it too, for the same underlying reason: those points are exactly
where a coarse boundary check happens to line up regardless of the sub-pixel shift. Both traps
were only caught by deliberately mutation-testing the passing test itself --- commenting out
the two \texttt{M41} / \texttt{M42} lines above and re-running the existing check suite, which
stayed fully green even with the real fix entirely removed, exposing the false-positive
``verified'' claim before it could ship. The fix was a dedicated check sampling a specific
internal texel boundary a rotated/flipped sprite's own blended edge produces --- a pixel
confirmed to shift from one exact, named RGB value to another under the same mutation, the
first check in the whole suite that actually moved when the offset was removed. The general
lesson generalizes well past D3D9 specifically: a test that passes both with and without the
code change it is supposed to be protecting is not testing anything, and the only reliable way
to find that out is to deliberately break the implementation and confirm the test notices.

\section{SpriteSortMode and SpriteEffects}

\cnaclass{SpriteSortMode} has five values: \texttt{Deferred}, \texttt{Immediate},
\texttt{Texture}, \texttt{BackToFront}, \texttt{FrontToBack}. Only the latter two actually
honor a sprite's \texttt{layerDepth} argument --- the others ignore it entirely, by design,
matching XNA's own documented sort-mode semantics.

\texttt{Texture} mode is worth a second look before moving on, because its sort key is a real,
documented gotcha rather than an implementation detail: \texttt{SpriteBatch}'s flush groups
draws by \texttt{std::stable\_sort(..., [](a, b) \{ return a.texture < b.texture; \})} ---
comparing raw \emph{texture pointers}. The contract this actually guarantees is only two
things, both verified directly by \texttt{SpriteBatchTests.cpp}'s
\texttt{TextureGroupsDrawsByTextureAndPreservesGroupOrder} case (Task~414): draws sharing a
texture end up adjacent, and within one texture's group, submission order is preserved
(\texttt{stable\_sort}, not \texttt{sort}). \emph{Which} texture's group renders first is
\textbf{not} part of the contract at all --- it depends on where the two \cnaclass{Texture2D}
objects happen to live in memory for that particular run, which can differ between runs of the
identical program. A game that submits sprites in scrambled A/B/A/B order and assumes
\texttt{Texture} mode will always draw all of \texttt{A} before all of \texttt{B} (or vice
versa) is relying on behavior the API never promised; only overlap-independent layouts (or an
explicit depth/layer separation) are safe under \texttt{Texture} mode. A dedicated pixel test
for the SDL\_RENDERER renderer (Task~668) exists specifically because a mock renderer can verify
grouping and order but not that the real draw dispatch binds the correct texture to each
reordered call --- the two are logically separate failure modes, and only a real renderer can
be wrong about the second one.

\cnaclass{SpriteEffects} is where this chapter's first concrete, still-open gap lives:
it is declared as a plain \texttt{enum class} --- \texttt{None=0}, \texttt{FlipHorizontally=1},
\texttt{FlipVertically=2} --- with \emph{no bitwise operators defined at all}. In real XNA,
\texttt{SpriteEffects} is a \texttt{[Flags]} enum, and combining both flags
(\texttt{FlipHorizontally | FlipVertically}, value 3) is a normal, supported call. In CNA,
reaching value 3 requires an explicit \texttt{static\_cast<SpriteEffects>(3)}, and --- more
tellingly --- \texttt{SpriteBatch::DrawString}'s own per-axis direction lookup tables (added
when the flip-mirroring bug below was fixed) deliberately only have three entries, indexed
0 through 2. Value 3 is functionally unreachable through the normal calling convention. This
is flagged directly in \texttt{docs/spritefont-support.md} \S7 as a genuine, minor,
still-open API-completeness gap, not a crash risk --- a combined horizontal-and-vertical
flip of drawn text simply is not expressible today.

\section{A real, shared bug: SpriteEffects flip only affected sampling, not layout}

Chapter~\ref{ch:design-philosophy} promised that this book would name specific, checkable
bugs rather than paraphrase a changelog, and this one is worth naming in full because it
affected every renderer at once, which is unusual --- most of the bugs catalogued in
Part~\ref{part:renderers}
are renderer-specific. Before Task~694, a \texttt{SpriteEffects} flip on
\texttt{DrawString} only flipped each individual glyph's own texture sampling in place; it
never mirrored the \emph{sequence or position} of the glyphs themselves, which is what FNA's
real implementation does via a pair of axis-direction/axis-mirrored terms. The fix lives in
the shared, renderer-agnostic \texttt{SpriteBatch.cpp} --- meaning it was one fix that
corrected the behavior identically on every one of CNA's graphics renderers simultaneously,
rather than something the chapters in Part~\ref{part:renderers} need to revisit renderer by renderer.

\section{A real, lasting side effect: Begin() changes GraphicsDevice.BlendState}
\label{sec:spritebatch-blendstate-leak}

The \texttt{BlendState} passed to \texttt{Begin()} (or the default \cnaclass{BlendState::AlphaBlend}
when omitted) is not scoped to that \texttt{Begin()} / \texttt{End()} pair --- it genuinely
overwrites \texttt{GraphicsDevice.BlendState} as a lasting side effect, exactly as real FNA's
own \texttt{SpriteBatch.cs} does, and \texttt{End()} does \emph{not} restore whatever
\texttt{BlendState} was active beforehand. A 3D draw call issued right after a
\texttt{SpriteBatch} pass, without the game explicitly reassigning \texttt{BlendState} itself,
silently inherits the \texttt{SpriteBatch} pass's blend mode:

\begin{lstlisting}
spriteBatch_->Begin(SpriteSortMode::Deferred, BlendState::Additive);
spriteBatch_->Draw(cornerTexture, Rectangle(0, 0, 4, 4), Color::White);
spriteBatch_->End();

// No BlendState reassignment here -- this quad draws with Additive, not AlphaBlend or Opaque,
// because SpriteBatch::Begin() changed GraphicsDevice.BlendState and End() never restores it.
device.DrawUserPrimitives(PrimitiveType::TriangleList, verts, 0, 2);
\end{lstlisting}

This was, until Task~956, a genuine bug rather than the faithfully-preserved FNA behavior
described above: \texttt{EasyGLSpriteBatchRenderer::Begin()} used to hardcode its own
\texttt{SrcAlpha} / \texttt{OneMinusSrcAlpha} blend factors on the real GL state regardless of
which \texttt{BlendState} was actually requested, and leave that hardcoded value in effect
after \texttt{End()} --- so a following 3D draw inherited leftover \emph{implementation}
garbage, not the requested, XNA-faithful blend mode. The fix (verified by a dedicated pixel
test that clears to gray, additive-blends a small corner sprite, then draws a full-viewport
opaque red 3D quad with no explicit blend state of its own, and checks that the quad's green
channel at a point away from the corner reads $\approx 100$ rather than $0$) makes
\texttt{Begin()} apply the real requested state and leave it there, matching FNA exactly rather
than merely avoiding a crash.

A second, related caveat is still open rather than fixed: on the EasyGL renderer, a
\textbf{full-backbuffer} \texttt{SpriteBatch} draw issued before any 3D draw in the same frame
breaks that frame's 3D rendering entirely (Task~933) --- investigated at length (window-title
readback state, viewport-reset ordering, and \texttt{RenderTarget2D} binding have each been
ruled out as the trigger) but not yet root-caused as of this writing. The blend-state-leak fix
above is a different bug with a different, narrower trigger (a \emph{non-full-backbuffer}
sprite specifically avoids Task~933's condition, which is exactly why the Task~956 regression
test draws its verification sprite into only a $4\times4$ corner) --- do not conflate the two
if you hit unexplained 3D-rendering corruption after a \texttt{SpriteBatch} pass on this
renderer; check whether the preceding sprite pass covered the full backbuffer before assuming
the blend-state fix above has regressed.

\section{SpriteFont}
\label{sec:spritefont-measure}

\cnaclass{SpriteFont} is deliberately not a \cnaclass{GraphicsResource} --- a texture atlas,
glyph bounds, cropping data, a character set, line spacing, character spacing, kerning data,
and an optional default character. Its constructor is \texttt{CNAEXT} and publicly callable,
a direct consequence of Chapter~\ref{ch:content-and-assets}'s content model: real XNA's
\cnaclass{SpriteFont} constructor is \texttt{internal}, reachable only from the compiled
\texttt{.xnb} content pipeline's own \texttt{SpriteFontReader}; CNA has no such reader, so the
constructor has to be public for anything to construct a font at all.

\texttt{MeasureString} exists for both \texttt{string} and \texttt{StringBuilder} --- the
\texttt{StringBuilder} overload was entirely missing until Task~421/423, and is now a
one-line forward to the \texttt{string} implementation, a simplification available in
C\texttt{++} precisely because there is no garbage-collection-pressure argument (the reason
FNA's own C\# implementation duplicates the algorithm instead of forwarding) against sharing
code here. A faithfully-preserved, slightly surprising detail: a trailing \texttt{"\textbackslash n"}
adds a full second empty line's worth of height to the measured result ---
\texttt{Y = 2 \(\times\) LineSpacing}, not \texttt{1 \(\times\)} --- because both the newline
handler and the loop's own unconditional final height-add fire. This matches real XNA's
behavior exactly and is preserved rather than ``fixed,'' consistent with
Chapter~\ref{ch:design-philosophy}'s fidelity-over-taste rule.

Character encoding deserves its own note: CNA's \texttt{String} is UTF-8
(\texttt{std::string}), and \cnaclass{SpriteFont} / \cnaclass{SpriteBatch} decode each glyph via
an internal \texttt{DecodeUtf8CodePoint} helper. Because \texttt{charcs} (the glyph-key type,
per Chapter~\ref{ch:design-philosophy}'s type-alias table) is 16 bits wide, only Basic
Multilingual Plane code points can serve as glyph keys at all; an invalid or truncated UTF-8
sequence decodes to \texttt{'?'} rather than throwing, and the byte index always advances by
at least one, so a malformed string can never cause an infinite decode loop.

A worked example ties \texttt{MeasureString} directly to the simplest \texttt{DrawString}
overload for the single most common real use --- centering a line of HUD text horizontally
around a screen-space point, which needs the string's own measured width before it can compute
where to start drawing:

\begin{lstlisting}
const std::string text = "Game Over";
Vector2 size = font.MeasureString(text);           // Vector2(width, height) in pixels
Vector2 centered(screenCenterX - size.X / 2.0f, screenCenterY - size.Y / 2.0f);

spriteBatch_->DrawString(font, text, centered, Color::White);
\end{lstlisting}

\section{What SpriteFont's support matrix actually says}

\texttt{docs/spritefont-support.md} verifies the property surface, \texttt{MeasureString},
single- and multi-glyph placement, newline handling, default-character fallback, single-axis
flip, and rotation/scale as pixel-correct on both \textbf{SDL\_Renderer} and \textbf{EasyGL}.
On \textbf{Vulkan} and \textbf{BGFX}, the underlying logic is the same shared C\texttt{++} ---
very likely correct by construction --- but has not been independently re-verified
specifically for \cnaclass{SpriteFont} (flagged again in
\texttt{docs/\allowbreak{}graphics-\allowbreak{}renderer-\allowbreak{}feature-\allowbreak{}matrix.md},
Task~861). Combined-axis flip is not
representable on any renderer at all, for the \cnaclass{SpriteEffects} reason above.

\section{Two real, fixed SDL\_RENDERER bugs, previewing the renderer-audit method}
\label{sec:sdl-renderer-two-bugs}

Two of the fifteen real bugs found during the SDL\_RENDERER renderer's own Phase~70 audit
(Chapter~\ref{ch:sdl-renderer} covers the full audit) live directly in \cnaclass{SpriteBatch}
rendering and are worth previewing here as a template for how this book reports a bug
throughout Part~\ref{part:renderers}: name the symptom, the root cause, and the fix. First,
\texttt{SDL\_RenderTextureRotated}'s pivot convention differed from XNA's own rotation-origin
convention; the fix offsets the destination rectangle rather than changing the rotation math
itself. Second --- more serious --- the \texttt{transformMatrix} argument to
\texttt{Begin()} was silently ignored \emph{entirely} on this renderer; any code passing a
non-identity transform matrix (a common pattern for camera-style 2D scrolling) simply drew as
if it had passed identity. The fix added a new rendering path using
\texttt{SDL\_RenderTextureAffine}, used only when the transform actually differs from
identity, so the common (and cheaper) untransformed path is unaffected.

\section{How a batch actually reaches the GPU: flush timing, not real GPU batching}
\label{sec:spritebatch-flush-timing}

The class is named \cnaclass{SpriteBatch}, and \texttt{SpriteSortMode::Deferred} --- the
default --- sounds like it defers submission to build one large GPU-side batch. Reading
\texttt{SpriteBatch.cpp}'s own \texttt{pushSprite} / \texttt{flushSingle} / \texttt{flushBatch}
directly shows a narrower reality worth being precise about, since it has real performance
implications for a porting engineer used to a true GPU-batched 2D renderer:

\begin{lstlisting}
void SpriteBatch::pushSprite(...)
{
    // ...
    if (sortMode_ == SpriteSortMode::Immediate)
        flushSingle(info);              // dispatched to the renderer right here, right now
    else
        spriteQueue_.push_back(info);   // deferred -- but only the SORT, not the draw call
}

void SpriteBatch::flushBatch()
{
    // ... optional std::stable_sort by layerDepth or texture pointer, per sort mode ...
    for (const SpriteInfo& s : spriteQueue_)
        flushSingle(s);                 // still one renderer->Draw() call per sprite
    spriteQueue_.clear();
}
\end{lstlisting}

Two real, checkable facts fall out of reading this directly rather than assuming from the
class name. First, \texttt{SpriteSortMode::Immediate} is not merely ``skip the sort'' ---
\texttt{pushSprite} routes straight to \texttt{flushSingle} for every \texttt{Draw()} call,
meaning each sprite reaches the renderer the instant it is submitted, with no queue involved at
all; this is the only sort mode where a game can safely change render state (a different
\cnaclass{Texture2D}'s \cnaclass{SamplerState}, say, via a raw renderer call) \emph{between}
two \texttt{Draw()} calls in the same \texttt{Begin()} / \texttt{End()} pair and have it take
effect exactly where expected, since nothing is reordered or delayed. Second, and more
consequential for the other four sort modes: ``batch'' here describes when the \emph{sort}
happens, not how the GPU submission happens. Every sprite, in every sort mode, still becomes
its own individual \texttt{renderer->Draw()} call inside the \texttt{flushBatch()} loop above
--- there is no vertex-buffer-building step that merges many sprites into one draw call the
way a true GPU-batched 2D renderer (or, within CNA's own codebase, a single
\cnaclass{VertexBuffer} submitted once via \texttt{DrawIndexedPrimitives}) would. A porting
engineer coming from an engine where ``sprite batching'' specifically means draw-call
coalescing should recalibrate: CNA's \cnaclass{SpriteBatch} coalesces \emph{sort decisions},
not draw calls, and a scene with a thousand sprites issues on the order of a thousand renderer
draw calls regardless of which non-\texttt{Immediate} sort mode is chosen.

A second, smaller finding lives in \texttt{pushSprite} itself: a
\texttt{System::ObjectDisposedException::ThrowIf} guard on the texture argument (Task~717).
Real FNA's own \texttt{SpriteBatch.Draw} does not guard against a disposed texture at all ---
its managed runtime simply throws whatever exception a stale reference happens to produce.
C\texttt{++} has no equivalent safety net: reaching \texttt{flushSingle} with a fully-disposed
\cnaclass{Texture2D} (its last owning \texttt{shared\_ptr} released, renderer pointer now
\texttt{nullptr}) would dereference a null \cnaclass{ITextureRenderer} reference through
\texttt{GetRenderer()} --- a guaranteed crash, not a graceful failure. The explicit guard is
real, deliberate hardening beyond parity, not a faithfully-reproduced FNA behavior.

\section{A parallel to the BlendState leak: DepthStencilState's own Task~803 history}

\S\ref{sec:spritebatch-blendstate-leak} above covers \texttt{BlendState} leaking \emph{out} of
a \texttt{SpriteBatch} pass --- \texttt{Begin()} sets it, \texttt{End()} never restores it,
matching real FNA exactly. \texttt{DepthStencilState} had the mirror-image bug, now fixed
(Task~803): reading \texttt{Begin()}'s own implementation shows the
\texttt{depthStencilState} parameter was, for a real stretch of this project's history,
accepted but never actually applied to \texttt{graphicsDevice\_} at all --- a
\texttt{SpriteBatch} pass silently inherited whatever \texttt{DepthStencilState} the game's
most recent 3D rendering (or each renderer's own construction-time default) had last left in
place, rather than FNA's real, documented default of \texttt{DepthStencilState::None} when the
caller passes \texttt{nullptr}. The two bugs are opposite shapes of the identical underlying
discipline failure --- one state object leaking \emph{into} a pass that should have reset it,
the other leaking \emph{out of} a pass that XNA specifies should not restore it --- and the fix
for both ended up at the same place: \texttt{Begin()} now always (re-)applies the resolved
\texttt{DepthStencilState} (\texttt{nullptr} $\to$ \texttt{DepthStencilState::None}, exactly
mirroring the existing \texttt{samplerState} $\to$ \texttt{SamplerState::LinearClamp}
resolution immediately below it in the same function), never leaving a previous
\texttt{Begin()}'s value in effect. A porting engineer relying on either state object's value
surviving across a \texttt{SpriteBatch} pass in either direction is relying on real,
verified-current XNA/FNA-faithful behavior for \texttt{BlendState}, and real, now-fixed
CNA-specific correctness for \texttt{DepthStencilState} --- not a coincidence of implementation
detail in either case.

\section{A worked example: the real per-glyph advance math behind DrawString}
\label{sec:spritebatch-drawstring-advance}

\S\ref{sec:spritefont-measure} above already named the raw ingredients --- \texttt{kerning\_},
a \cnaclass{Vector3} per glyph packing \texttt{(leftBearing, width, rightBearing)} --- but
\texttt{MeasureString} and \texttt{DrawString} are two independent implementations of the
identical walk, in two different files (\texttt{SpriteFont.cpp} and \texttt{SpriteBatch.cpp}
respectively), and it is worth confirming directly, rather than assuming, that they agree.
Reading \texttt{SpriteBatch::DrawString} shows the same three-line shape, glyph by glyph:

\begin{lstlisting}
const Vector3& cKern = spriteFont.kerning_[index];
if (firstInLine)
{
    curOffset.X += std::abs(cKern.X);   // first glyph: ABS of left bearing
    firstInLine = false;
}
else
{
    curOffset.X += spacing_ + cKern.X;  // later glyphs: inter-char spacing + left bearing, no abs
}
curOffset.X += cKern.Y + cKern.Z;       // every glyph: width + right bearing
\end{lstlisting}

The \texttt{std::abs()} on the very first glyph's left bearing --- and only the first --- is the
detail easiest to miss reading this quickly, and it is not a typo: a font whose first glyph
happens to have a \emph{negative} left bearing (the glyph's own ink starts left of its nominal
advance origin, a real, common shape for italic or swash characters) must not have that negative
value shift the whole string's start position backward off of \texttt{position}, but a
\emph{later} glyph's negative left bearing legitimately tucks it closer to the previous
character --- exactly the kerning behavior italic fonts rely on. Confirming \texttt{MeasureString}
(\S\ref{sec:spritefont-measure}) applies the identical \texttt{abs()}-only-for-the-first-glyph
rule is what makes a caller's own centering math (measure, then draw at the measured offset)
trustworthy: if the two implementations disagreed on this one case, text would measure as one
width and render at another.

\texttt{modules/renderers/easygl/examples/\allowbreak{}easygl\_spritefont\_multiglyph\_spacing\_test.cpp} (Task~425) is a real,
minimal fixture built specifically to prove the non-first-glyph branch is genuinely exercised,
not merely present in the source. Its two-glyph font --- \texttt{'A'} and \texttt{'B'}, both
$8\times8$, both with zero-valued kerning \texttt{(0, 8, 0)} --- and a \texttt{spacing} of
\texttt{4.0f} let the whole placement be hand-verified in whole pixels:

\begin{lstlisting}
// 'A' (first in line): curOffset.X = abs(0) = 0        -> glyph drawn at dest x = 2 + 0 = 2
// after 'A':           curOffset.X += 8 + 0 = 8         (width + rightBearing)
// 'B' (not first):     curOffset.X += 4 + 0 = 12        -> glyph drawn at dest x = 2 + 12 = 14
\end{lstlisting}

Drawing \texttt{"AB"} at \texttt{position = (2, 2)} therefore places \texttt{'A'} at screen
$x \in [2, 10)$ and \texttt{'B'} at $x \in [14, 22)$ --- an exact 4-pixel gap at $x \in [10, 14)$
that must stay background-colored for the test to pass. The real test's own check set is worth
noting for what it rules out beyond simple placement: sampling a point \emph{inside} the gap
confirms \texttt{spacing} was genuinely applied (not silently zero, and not some other stray
value), while sampling each glyph's own distinct fill color (white for \texttt{'A'}, green for
\texttt{'B'}) at its expected position rules out an index mix-up between the two glyphs ---
a bug shape that a single-glyph test, no matter how many pixels it samples, structurally cannot
detect at all, since there would be no second glyph to place at the wrong index.

\section{A real exception-type gap: nested \texttt{Begin()} throws the wrong type entirely}
\label{sec:spritebatch-wrong-exception-type}

Calling \texttt{Begin()} a second time without an intervening \texttt{End()} --- or calling
\texttt{End()} before any \texttt{Begin()} at all --- is a real, commonly-hit misuse case (a
forgotten \texttt{End()} inside an early-return branch, most often), and CNA does throw for it.
Reading the real implementation shows exactly what it throws, though, and it is not the type a
reader would reasonably expect from real XNA:

\begin{lstlisting}
if (begun)
    throw std::runtime_error("Begin has been called before calling End.");
// ...
if (!begun)
    throw std::runtime_error("End was called, but Begin has not yet been called.");
\end{lstlisting}

Real XNA's own \texttt{SpriteBatch.Begin()} throws \texttt{InvalidOperationException} for the
identical misuse, and \cnans{System}'s own \cnaclass{InvalidOperationException} type exists in
this codebase, real and used elsewhere (Chapter~\ref{ch:state-objects} and others cite it
directly) --- but \cnaclass{SpriteBatch} itself throws a plain \texttt{std::runtime\_error}
here instead, confirmed by reading both guard sites directly rather than assumed from the
message text alone. The two exception hierarchies do not meet until \texttt{std::exception}
itself: \cnaclass{InvalidOperationException} derives from \cnaclass{SystemException}, which
derives from \cnaclass{System::Exception} --- a hierarchy entirely separate from
\texttt{std::runtime\_error}'s own \texttt{std::exception} lineage, confirmed by reading
\texttt{InvalidOperationException.hpp} / \texttt{SystemException.hpp} directly.

\begin{lstlisting}
// Worked example: a porting engineer's reasonable, XNA-faithful guard --
// this catch block never fires against SpriteBatch's own Begin()/End()
// misuse, even though it would fire against the identical misuse on many
// other CNA classes that do throw the real InvalidOperationException.
try
{
    spriteBatch_->Begin();
    spriteBatch_->Begin();  // forgot the End() above -- a real, easy mistake
}
catch (const System::InvalidOperationException& ex)
{
    // Never reached: SpriteBatch::Begin() throws std::runtime_error here,
    // which this catch clause's type does not match at all.
    LogAndRecoverGracefully(ex);
}
// The exception propagates past this handler entirely, caught only by a
// broader catch (const std::exception&) further up the call stack, if one
// exists -- or terminates the program if it does not.
\end{lstlisting}

This is a narrower, more precisely scoped gap than ``wrong error message'' or ``doesn't throw at
all'': the throw itself is correct and reliable (a real, hittable misuse genuinely does
terminate the batch and report a real diagnostic message), and any code written against a
generic \texttt{catch (const std::exception\&)} handler --- the common, safe default in
ordinary C\texttt{++} --- catches it exactly as expected, message text and all. The gap only
bites a porting engineer who specifically wrote a typed \cnaclass{InvalidOperationException}
catch clause around \cnaclass{SpriteBatch} usage, reasonably expecting the same exception
hierarchy this book's other chapters confirm many other CNA classes do use faithfully for the
identical real-XNA exception type.

\section{Where this leaves a porting engineer}

\cnaclass{SpriteBatch} is, by a wide margin, the most exhaustively cross-renderer-verified
class in CNA --- every renderer chapter in Part~\ref{part:renderers} returns to it, because every renderer has
to implement it, and because its bugs (an ignored transform matrix, a mis-scoped flip, a
silently-downgraded sampler filter) are exactly the class of bug that only surfaces under
real pixel-readback testing rather than ``it compiles and something appears on screen.'' If
you are porting existing 2D XNA/FNA game code (this book's migration chapter), this is also
the API you can trust most: it is the one every renderer is held to the same standard on.
